% Part: first-order-logic
% Chapter: beyond
% Section: higher-order-logic

\documentclass[../../../include/open-logic-section]{subfiles}

\begin{document}

\olfileid[pt]{fol}{byd}{hol}

\olsection{Lógica de Ordem Superior}

Passar da lógica de primeira ordem para a de segunda ordem nos permitiu
falar sobre conjuntos de objetos no !!{domain} de primeira ordem, dentro da
linguagem formal. Por que parar aí? Por exemplo, a lógica de terceira
ordem deveria nos permitir lidar com conjuntos de conjuntos de objetos,
ou talvez até conjuntos que contenham tanto objetos quanto conjuntos de
objetos. E a lógica de quarta ordem nos permitirá falar sobre conjuntos
de objetos desse tipo. Como se pode imaginar, pode-se iterar essa ideia
arbitrariamente.

Na prática, a lógica de ordem superior é frequentemente formulada em
termos de funções em vez de relações. (Módulo as identificações naturais,
essa diferença é inessencial.) Dadas algumas ``ordenações'' básicas $A$, $B$,
$C$,~\dots (que agora chamaremos de ``tipos''), podemos criar novos
estipulando
\begin{quote}
Se $\sigma$ e $\tau$ são tipos finitos, então $\sigma \to \tau$ também o é.
\end{quote}
Pense nos tipos como ``rótulos'' sintáticos que classificam os objetos que
queremos em nosso !!{domain}; $\sigma \to \tau$ descreve aqueles objetos que
são funções que levam objetos de tipo~$\sigma$ a objetos de
tipo~$\tau$. Por exemplo, podemos querer ter um tipo $\Omega$ de
valores de verdade, ``verdadeiro'' e ``falso'', e um tipo $\Nat$ de números
naturais. Nesse caso, pode-se pensar em objetos de tipo $\Nat \to
\Omega$ como relações unárias, ou subconjuntos de $\Nat$; objetos de tipo
$\Nat \to \Nat$ são funções de números naturais para números naturais;
e objetos de tipo $(\Nat \to \Nat) \to \Nat$ são ``funcionais'',
isto é, funções de tipo superior que levam funções a números.

Como no caso da lógica de segunda ordem, pode-se pensar na lógica de
ordem superior como uma espécie de lógica de múltiplas ordenações, em que
há uma ordenação para cada tipo de objeto que queremos considerar. Mas é
usualmente mais
claro simplesmente definir a sintaxe da lógica de tipos superiores a
partir dos fundamentos. Por exemplo, podemos definir um conjunto de tipos
finitos indutivamente, da seguinte forma:
\begin{enumerate}
\item $\Nat$ é um tipo finito.
\item Se $\sigma$ e $\tau$ são tipos finitos, então $\sigma
  \to \tau$ também o é.
\item Se $\sigma$ e $\tau$ são tipos finitos, então $\sigma \times
  \tau$ também o é.
\end{enumerate}
Intuitivamente, $\Nat$ denota o tipo dos números naturais,
$\sigma \to \tau$ denota o tipo de funções de $\sigma$ para
$\tau$, e $\sigma \times \tau$ denota o tipo de pares de objetos,
um de $\sigma$ e outro de $\tau$. Podemos então definir um conjunto de
termos indutivamente, da seguinte forma:
\begin{enumerate}
\item Para cada tipo $\sigma$, há um estoque de variáveis $x$, $y$,
  $z$, \dots de tipo $\sigma$
\item $\Obj 0$ é um termo de tipo $\Nat$
\item $\Obj S$ (sucessor) é um termo de tipo $\Nat \to \Nat$
\item Se $s$ é um termo de tipo $\sigma$, e $t$ é um termo de tipo
  $\Nat \to (\sigma \to \sigma)$, então $\Obj{R}_{st}$ é um termo de tipo
  $\Nat \to \sigma$
\item Se $s$ é um termo de tipo $\tau \to \sigma$ e $t$ é um
  termo de tipo~$\tau$, então $s(t)$ é um termo de tipo $\sigma$
\item Se $s$ é um termo de tipo~$\sigma$ e $x$ é uma variável de
  tipo~$\tau$, então $\lambd[x][s]$ é um termo de tipo $\tau \to \sigma$.
\item Se $s$ é um termo de tipo~$\sigma$ e $t$ é um termo de
  tipo~$\tau$, então $\tuple{s, t}$ é um termo de tipo $\sigma \times
  \tau$.
\item Se $s$ é um termo de tipo~$\sigma \times \tau$, então $p_1(s)$ é um
  termo de tipo~$\sigma$ e $p_2(s)$ é um termo de tipo~$\tau$.
\end{enumerate}
Intuitivamente, $\Obj{R}_{st}$ denota a função definida recursivamente por
\begin{align*}
\Obj{R}_{st}(0) & = s \\
\Obj{R}_{st}(x+1) & = t(x, R_{st}(x)),
\end{align*}
$\tuple{s, t}$ denota o par cuja primeira componente é~$s$ e cuja
segunda componente é~$t$, e $p_1(s)$ e~$p_2(s)$ denotam o primeiro
e o segundo elementos (``projeções'') de~$s$. Finalmente, $\lambd[x][s]$
denota a função~$f$ definida por
\[
f(x) = s
\]
para qualquer~$x$ de tipo~$\sigma$; assim o item (6) nos dá uma forma de
compreensão, permitindo-nos definir funções usando
termos. As !!{formula}s são construídas a partir de afirmações de
!!{identity} $\eq[s][t]$ entre termos do mesmo tipo, dos conectivos
proposicionais usuais e da quantificação de ordem superior. Pode-se então
tomar os axiomas do sistema como as equações básicas que governam os
termos definidos acima, juntamente com as regras usuais da lógica com
quantificadores e !!{identity}.

Se se aumenta o sistema de tipos finitos com um tipo $\Omega$ de valores
de verdade, é preciso incluir axiomas que governem seu uso também. De
fato, se se for astuto, pode-se eliminar !!{formula}s complexas
inteiramente, substituindo-as por termos de tipo $\Omega$!{} O sistema de
prova pode então ser modificado de acordo. O resultado é essencialmente a
\emph{teoria simples dos tipos} apresentada por Alonzo Church na década de
1930.

Como no caso da lógica de segunda ordem, há diferentes versões de
semântica de tipos superiores que se pode querer usar. Na versão completa,
variáveis de tipo $\sigma \to \tau$ variam sobre o conjunto de \emph{todas}
as funções dos objetos de tipo~$\sigma$ para objetos de tipo~$\tau$.
Como se pode imaginar, essa semântica é forte demais para admitir um
sistema de !!{derivation} completo e efetivo. Mas pode-se considerar uma
semântica mais fraca, em que !!a{structure} consiste em conjuntos de
elementos $T_\tau$ para cada tipo $\tau$, juntamente com operações
apropriadas para aplicação, projeção, etc. Se os detalhes forem
desenvolvidos corretamente, pode-se obter teoremas de completude para os
tipos de sistemas de !!{derivation} descritos acima.

A lógica de ordem superior é atraente porque fornece um arcabouço no
qual podemos embutir boa parte da matemática de maneira natural: a partir
de $\Nat$, pode-se definir números reais, funções contínuas e assim por
diante. É também particularmente atraente no contexto da lógica
intuicionista, pois os tipos têm interpretações ``construtivas'' claras.
De fato, pode-se desenvolver versões construtivas da semântica de tipos
superiores (baseadas na lógica intuicionista, em vez da clássica) que
esclarecem bastante essas interpretações construtivas e são, em muitos
aspectos, mais interessantes do que as contrapartes clássicas.

\end{document}
